% Source-aligned English for printed p.154.
\SourcePage{154}

\LayerTarget{tohoku:def:presheaf-and-sheaf:p2}
\noindent that the $F(U)$ are groups, rings, and so on, and that the
$\varphi_{VU}$ are homomorphisms, one obtains the notion of a presheaf or
sheaf of groups, rings, and so on. More generally, one can define a presheaf
or sheaf with values in a given category; see 1.1. Presheaves or sheaves on
$X$ with values in a given category themselves form a category, with
morphisms defined as morphisms of inductive systems. Presheaves or sheaves on
$X$ with values in an additive category, for example the category of abelian
groups, form an additive category; presheaves and sheaves of abelian groups
even form abelian categories. For short, we shall say \textit{abelian sheaf}
and \textit{abelian presheaf} for a sheaf or presheaf of abelian groups.

\LayerTarget{tohoku:remark:sheaves-vs-presheaves-abelian}
Note, however, that the identity functor sending an abelian sheaf to its
underlying abelian presheaf is left exact but not exact. Given a homomorphism
of sheaves $u:F\longrightarrow G$, its cokernel as a homomorphism of
presheaves is
\[
Q(U)=G(U)/\operatorname{Im}F(U),
\]
which in general is not a sheaf. Its cokernel as a homomorphism of sheaves is
the sheaf associated to the presheaf $Q$, as described below. We shall not
dwell on these presently well-known matters; see [4] and Godement's book [9].

\LayerTarget{tohoku:def:etale-space-from-presheaf}
Let $F$ be a presheaf of sets on $X$. For every $x\in X$, put
\[
F(x)=\varinjlim_{U\ni x}F(U),
\]
where the direct limit is taken over the filtered ordering of the open
neighborhoods $U$ of $x$. On the set $\overline F$, the union of the $F(x)$,
put the topology generated by the subsets of the form $A(f)$: for every open
$U\subset X$ and $f\in F(U)$, $A(f)$ is the set of canonical images $f(x)$ of
$f$ in $F(x)$ for $x\in U$. With this topology, the natural map
$\overline F\longrightarrow X$ is a local homeomorphism. Following Godement,
we call $\overline F$ an \textit{étalé space} over $X$.

\LayerTarget{tohoku:prop:sheaves-etale-spaces-equivalence}
Conversely, an étalé space $E$ over $X$ naturally defines a sheaf
$\mathfrak F(E)$, which assigns to an open $U$ the set of continuous sections
of $E$ over $U$. The étalé space $\overline{\mathfrak F(E)}$ associated to
$\mathfrak F(E)$ is canonically identified with $E$. Starting from a
presheaf $F$ on $X$, there is also a canonical homomorphism
\[
F\longrightarrow\mathfrak F(\overline F),
\]
since each $f\in F(U)$ defines the continuous section
$x\longmapsto f(x)$ of $\overline F$ over $U$. This homomorphism is an
isomorphism if and only if $F$ is a sheaf. Thus: (i) \textit{the notion of a
sheaf of sets on $X$ is equivalent to that of an étalé space over $X$}; more
precisely, this gives an equivalence from the category of sheaves of sets on
$X$ to the category of étalé spaces over $X$.

\LayerTarget{tohoku:construction:sheafification:p1}
(ii) \textit{Every presheaf $F$ on $X$ has an associated sheaf
$\mathfrak F(\overline F)$ and a homomorphism
$F\longrightarrow\mathfrak F(\overline F)$, which is an isomorphism if and
only if $F$ is a sheaf.} Moreover, $\mathfrak F(\overline F)$ is a

\TohokuContinuation{155}

